% Generated draft: requires mathematical and source review.
\documentclass[11pt]{article}
\usepackage[margin=25mm]{geometry}
\usepackage{fontspec}
\setmainfont{Latin Modern Roman}
\usepackage{amsmath,amssymb,mathtools,unicode-math}
\setmathfont{Latin Modern Math}
\usepackage{enumitem,needspace,xurl,hyperref}
\setlength{\parindent}{0pt}
\setlength{\parskip}{5pt}
\setlength{\emergencystretch}{3em}
\newcommand{\N}{\mathbb{N}}
\newcommand{\Z}{\mathbb{Z}}
\newcommand{\Q}{\mathbb{Q}}
\newcommand{\R}{\mathbb{R}}
\newcommand{\C}{\mathbb{C}}
\newcommand{\F}{\mathbb{F}}
\newcommand{\A}{\mathbb{A}}
\newcommand{\PP}{\mathbb{P}}
\newcommand{\E}{\mathbb{E}}
\newcommand{\eps}{\varepsilon}
\newcommand{\dd}{\,\mathrm{d}}
\begin{document}
\section*{UnsolvedMath ID: 30004616}
\textbf{Source problem code:} catalogue entry
\subsection{Definitions and mathematical statement}
Let \(N\) be a positive integer, as required for a conductor or paramodular level. An abelian surface has dimension \(2\), and an abelian fourfold has dimension \(4\). Write \([X]_{\Q}\) for the \(\Q\)-isogeny class of an abelian variety \(X/\Q\), and let \(\operatorname{cond}(X)\) denote its arithmetic conductor. Define
\[
A_N=\left\{[A]_{\Q}: \dim A=2,\ \operatorname{cond}(A)=N,\ \operatorname{End}_{\Q}(A)=\Z\right\}.
\]
Here \(\operatorname{End}_{\Q}(A)\) is the ring of endomorphisms defined over \(\Q\). The source also specifies
\[
B_N=\left\{[B]_{\Q}: \dim B=4,\ \operatorname{cond}(B)=N^2,\ B\text{ has quaternionic multiplication}\right\},
\]
but does not specify the precise quaternionic-multiplication condition. In particular, it does not identify the quaternion algebra, the field of definition of its action, or additional endomorphism or simplicity conditions.

Let \(P_N\) denote the source's cuspidal nonlift Siegel paramodular newforms of level \(N\) with rational Hecke eigenvalues, modulo multiplication by a nonzero scalar. The original statement further qualifies these forms as ``suitable.'' The weight, the meaning of this suitability qualification, and the precise lift and newform conventions are not supplied, so this description does not yet determine a precise set \(P_N\).

Subject to resolving these qualifications, the proposed correspondence asserts that for every positive integer \(N\) there exists a bijection
\[
\Phi_N:A_N\cup B_N\longrightarrow P_N
\]
such that
\[
\begin{aligned}
L_{\mathrm{HW}}(A,s)&=L_{\mathrm{spin}}(\Phi_N([A]_{\Q}),s)
&&\text{for every }[A]_{\Q}\in A_N,\\
L_{\mathrm{HW}}(B,s)&=\bigl(L_{\mathrm{spin}}(\Phi_N([B]_{\Q}),s)\bigr)^2
&&\text{for every }[B]_{\Q}\in B_N.
\end{aligned}
\]
Here \(L_{\mathrm{HW}}\) denotes the Hasse--Weil \(L\)-function of the abelian variety, and \(L_{\mathrm{spin}}\) denotes the spin \(L\)-function of the form. The source does not specify their normalization, completion, or conventions at primes dividing \(N\). These choices must be fixed compatibly before the displayed equalities are fully precise. This is the supplied conjectural target, with unresolved hypotheses explicitly retained.
\subsection{Short English statement}
For each positive level \(N\), the proposed correspondence pairs the specified rational abelian surfaces and quaternionic-multiplication fourfolds with rational-eigenvalue paramodular forms. The associated Hasse--Weil \(L\)-function should equal the spin \(L\)-function for a surface and its square for a fourfold. The fourfold conditions and the class of eligible forms require clarification.
\subsection{Sources}
\begin{itemize}
\item Supplied local UnsolvedMath registry record, \texttt{data/unsolvedmath/problems.json}, permanent ID \(30004616\), source problem code \texttt{catalogue entry}, entitled ``Paramodular Correspondence for Abelian Surfaces and Fourfolds.'' The supplied record attributes the conjecture to Brumer--Kramer and names no individual proposer.
\item Supplied report citation: \emph{Moduli spaces and Modular forms} (2022), Oberwolfach Reports, pp.~301--359, DOI \texttt{10.4171/owr/2021/6}. \url{https://doi.org/10.4171/owr/2021/6}. The background describes Oberwolfach Reports as records of mathematical workshops and their open-problem discussions.
\item Supplied literature reference: Brumer--Kramer, ``paramodular conjecture.'' \url{https://arxiv.org/abs/1004.4699}.
\item Supplied literature reference: Boxer--Calegari--Gee--Pilloni, ``modularity theorems.'' \url{https://arxiv.org/abs/2502.20645}.
\item Supplied literature reference: ``Oberwolfach report.'' \url{https://doi.org/10.4171/owr/2021/6}.
\item Literature assessment dated exactly \texttt{2026-08-22}: the supplied curation labels the problem open, identifies the proposed surface and quaternionic-multiplication fourfold correspondence as the Brumer--Kramer paramodular conjecture, and states that it is not established for all levels and forms. This is supplied dated curation, not an independently verified assessment; the background expressly requests independent verification of time-sensitive status and cited claims.
\item The registry labels the statement \texttt{corrected\_verified} and reports that it was checked against the official Oberwolfach report. These are supplied verification claims, not independent verification by this draft. The proposed year is recorded as 2021, while the supplied report citation uses 2022 and its DOI contains 2021; these particulars are retained without reconciliation.
\end{itemize}
\end{document}
